\section{Introduction}
A word response can have constant signed Hankel rank and constant
nonnegative rank at every cut, yet require increasing positive state
cardinality when its cut factorizations must use compatible stochastic
updates. Positive realization, rather than ordinary linear realization,
is therefore the relevant formalism. We study the controlled numerical
response
\begin{equation}\label{eq:target}
 \Pp(B=b\mid x,w,j)=\frac{1+b\rho\ip{e_j}{U_wx}}2,
 \qquad b\in\{-1,1\},
\end{equation}
for a finite orthogonal alphabet, a unit seed and one selected terminal
coordinate query. Our basic resource is nonuniform clocked atomic-row
positive-realization label width. The known horizon, row tables and
exact real arithmetic are free in that model. A separate theorem below
charges them in a uniform finite-bit command-stream model.

There are three contributions to the present treatment. First,
Section~\ref{sec:hankel} writes the exact finite-horizon object in the
language of normalized positive Hankel factorizations and compatible
future-response polytopes. This specializes the classical
shift-invariant-cone principle; neither positivity nor compatibility is
claimed as a new realization concept. For the planar numerical family,
ordinary rank, nonnegative rank and normalized positive rank all equal
three at every cut. This identifies precisely what separate
factorization omits.

Second, write $\Gamma_{\calA}(k,B)$ for optimized length-$B$ word
quantization loss and $\lambda_{\calA}(k,a)$ for the least common
dilation of a $k$-generator polytope containing $a$ times the unit ball.
Theorem~\ref{thm:rate44} proves the universal comparison
\begin{equation}\label{eq:intro-rate44}
 \sup_{B\ge1}\frac{-\log(1-\Gamma_{\calA}(k,B))}{B}
       \le \log\lambda_{\calA}(k,a).
\end{equation}
It uses arbitrarily many independent packets to remove the initial
inradius loss. Thus the paired profiles are related for every finite
orthogonal alphabet, including noncommuting alphabets. Equality and
optimality among all hidden realizations are not assumed.

Third, Theorem~\ref{thm:uniform44} gives a table-free finite-bit compiler
for finitely many planar angles having fixed linear-space digit
procedures. Its parameter search, dyadic probabilities, arithmetic
scratch space, horizon counter and finite random words require
$O(\log N+\log(1/\varepsilon))$ work bits. Every command is processed
once; the output error is uniformly at most a prescribed
$\varepsilon>0$. For the explicit badly approximable algebraic-angle
families the work-space order is exactly $\Theta(\log N)$, while the
clean stochastic register retains the arithmetic label exponent.
This is not an exact fair-bit implementation of arbitrary real rows,
and the sharp coefficient of work space is not determined.

The packet, interval and enclosure theorems of the preceding development
are reproduced with their proofs in Section~\ref{sec:profiles}.
The packing--dilation balance there is a conditional meta-corollary,
not a classification of finite alphabets. The arithmetic results below
are retained: all of their analytic proofs remain in the current paper,
with the branch of the finite-type minimum displayed explicitly.

Our arithmetic conclusions go beyond uniform bad approximation. For $\alpha\in\R^r$ with $1,\alpha_1,\ldots,\alpha_r$ rationally independent, write
\begin{equation}\label{eq:type-intro}
 \omega^*(\alpha)=\limsup_{\norm m_1\to\infty}
       \frac{-\log\norm{m\cdot\alpha}_{\R/\Z}}{\log\norm m_1}.
\end{equation}
This is the ordinary dual exponent $\omega_{1,r}$ of the row vector $\alpha$ in \cite[Definition 1]{BL06}, with an equivalent fixed-dimensional norm; it is not the uniform exponent $\widehat\omega_{1,r}$ or the simultaneous exponent $\omega_{r,1}$. The alphabet consists of the identity and the rotations through $2\pi\alpha_i$. Throughout, $r$, the alphabet, $0<\rho\le1/10$ and
$0\le\varepsilon<\rho/(2\sqrt2)$ are fixed.

\begin{theorem}[Minimal Diophantine type]\label{thm:metric-main}
If $\omega^*(\alpha)=r$, then
\begin{equation}\label{eq:log-law}
 \lim_{N\to\infty}\frac{\log W_{N,\varepsilon}(\alpha)}{\log N}
                         =\frac{r}{2r+1}.
\end{equation}
For Lebesgue-almost every $\alpha$, and every fixed $\delta>0$, there are positive constants $c,C$, depending on the fixed $\alpha,\delta,r,\rho,\varepsilon$, such that, for $N\ge2$,
\begin{equation}\label{eq:ae-log-bounds}
 c\frac{N^{r/(2r+1)}}{(\log(N+2))^{2r(1+\delta)/(2r+1)}}
 \le W_{N,\varepsilon}(\alpha)
 \le C N^{r/(2r+1)}(\log(N+2))^{2r^2(1+\delta)/(2r+1)}.
\end{equation}
The upper inequalities are achieved by exact common-row machines. No uniformity near rational relations or in growing $r$ is asserted.
\end{theorem}

The exponent in \eqref{eq:log-law} need not imply a two-sided constant-factor law. The earlier dual badly approximable theorem supplies the stronger $\Theta$ conclusion on its stated class, and is recovered in Corollary~\ref{cor:ba}. Theorem~\ref{thm:metric-main} instead allows arbitrarily small Diophantine constants along subsequences. The quantitative arithmetic ingredients are classical; the memory conclusions concern their compatible realization and all-hidden-state converse.

There is also a genuine failure of a single exponent, even for one fixed irrational angle.
\begin{theorem}[Irrational horizons and Liouville fluctuations]\label{thm:liouville-main}
For every irrational $\alpha$ there are horizons $M_n\to\infty$ such that
$W_{M_n,\varepsilon}(\alpha)=\Theta(M_n^{1/3})$. If $\alpha$ is Liouville, then
\begin{equation}\label{eq:liouville-summary}
 \liminf_{N\to\infty}\frac{\log W_{N,\varepsilon}(\alpha)}{\log N}=0,
 \qquad
 \frac13\le\limsup_{N\to\infty}\frac{\log W_{N,\varepsilon}(\alpha)}{\log N}
 \le\frac12.
\end{equation}
Nevertheless $W_{N,\varepsilon}(\alpha)\to\infty$. All displayed upper bounds have exact realizations.
\end{theorem}

The upper limit in \eqref{eq:liouville-summary} is not determined here. The result establishes oscillation between subpower subsequences and a positive-power subsequence; it does not replace that distinction by a guessed exponent for all irrationals.


\paragraph{Provenance of the ingredients.}
The stochastic Hankel normal form and its finite separation dual are
normalizations of established realization and convex-separation
arguments \cite{BF04,Ohta20,BPP15}. Their role is to place the resource,
not to claim priority for positive cones. The endpoint mechanism and
arithmetic laws are inherited \cite{GTF35,GTF42,GTF43}. The new statements
are the distortion-rate comparison and the uniform finite-bit
realization with its charged resource ledger. Dirichlet approximation,
continued fractions, metric tail estimates and the algebraic examples
remain classical. Section~\ref{sec:comparison} compares the precise
objectives and norms, rather than treating adjacent theories as absent.
